% ============================================================
% conclusion.tex — Conclusion and Future Directions
% ============================================================

This chapter draws together the thesis's central findings,
considers what they imply for securities litigation research
and reform evaluation, and identifies the most productive
directions for extending this work.

\section{Summary of Contributions}

This thesis asked how the PSLRA, court geography, and
observable case characteristics jointly shape the timing and
type of resolution in federal securities class actions.
Applying a competing-risks survival framework to 12,968 cases
from the FJC IDB 1990--2024, we show
that the answer depends materially on the statistical
framework chosen---and that the standard single-outcome
approach systematically distorts the picture of securities
litigation dynamics. Four substantive contributions follow from that reframing.

First, the PSLRA's association with case outcomes is
time-varying rather than uniform. The dismissal association
is concentrated early and fades over the litigation duration,
consistent with the statute's pleading-stage design. The
settlement evidence is less stable---it attenuates under
reweighting, widens under cluster-robust adjustment, and
varies across circuits. Together, the evidence points to an early dismissal filter rather than a uniform shift in litigation dynamics.

Second, conditioning on observable case-mix shifts via IPTW
leaves the dismissal association elevated but materially
attenuates the settlement association. These are sensitivity
comparisons rather than a formal decomposition---some
weighting covariates may themselves lie on post-PSLRA
pathways, limiting causal interpretation. The settlement
association also fails to survive cluster-robust adjustment
($p = 0.158$) and weakens in the two largest circuit-specific
sub-samples. Across these checks, the dismissal finding is
the more stable of the two.

Third, circuit geography emerges as a dominant structural
predictor of case trajectories. Relative to the Second
Circuit, every other major circuit is more settlement-prone,
and the frailty results suggest that settlement outcomes are
more sensitive than dismissal outcomes to local judicial
culture and bar practices.

Fourth, the MDL sign-reversal finding illustrates why
competing-risks analysis is especially valuable in this
setting. MDL cases resolve more slowly at the cause-specific
level but have a higher cumulative settlement probability once
the competing dismissal pathway is modeled directly.

\section{Future Directions}

The most tractable near-term steps are district-level
disaggregation within the existing IDB and narrow-window
re-estimation using the current covariate set. Merging with
CourtListener, Stanford, and CRSP records are more ambitious
data-integration projects that would require substantial
additional effort.

\subsection{Unmeasured Confounders and Data Enrichment}

The most consequential limitation we encountered is the
absence of case-specific merit signals from the FJC's
administrative data. Three categories of unmeasured
confounders warrant investigation in future work.

\paragraph{Judicial ideology and case-management norms.}
The frailty models detect residual settlement heterogeneity
at the circuit level ($\hat{\theta} = 0.130$) but cannot
identify its source. Judge identifiers recoverable from
CourtListener docket records, combined with Judicial Common
Space ideology scores \citep{epstein2007}, would allow for extensions to judge-level random effects. That would permit
direct estimation of whether individual judicial approaches
to case management, discovery scheduling, and motion practice
drive the geographic patterns we observe.

\paragraph{Plaintiff firm sophistication.}
Law firm identity is a strong predictor of litigation
outcomes \citep{wang2022}, yet the IDB lacks
attorney identifiers. CourtListener integration would
provide title-based class-action validation, lead-counsel
identification, and docket-activity proxies for litigation
complexity. These variables could be incorporated as
additional covariates or as stratification factors in the
competing-risks models, sharpening the interpretation of the
PSLRA's filtering association.

\paragraph{Economic conditions and market context.}
Securities fraud allegations are often triggered by sharp
stock-price declines, and the business cycle may influence
both the volume and the composition of filings. Merging with
CRSP market data would allow for examination of whether
economic conditions at the time of filing moderate the PSLRA
association and whether the post-PSLRA compositional shift
that our IPTW analysis adjusts for is partly driven by
market-cycle effects not captured by the current covariates.

\subsection{The Second Circuit Pattern}

The most distinctive feature of our geographic analysis is not that the
Second Circuit has the highest dismissal CIF, but that it combines the
lowest settlement baseline with persistently high dismissal incidence.
Every other major circuit's settlement hazard is 1.1 to 4.3 times that of
the Second Circuit, while the Fifth Circuit slightly exceeds it in
cumulative dismissal incidence. Potential mechanisms---elite defense
counsel, judicial expertise, and demanding motion-to-dismiss
standards---are plausible but cannot be disentangled with the current data.

Three extensions would sharpen this analysis. First,
disaggregating the Second Circuit into its constituent
districts (SDNY, EDNY, NDNY, WDNY, and the District of
Vermont) would reveal whether the pattern is driven primarily
by the Southern District of New York or is circuit-wide.
Second, comparing pre- and post-PSLRA trends within the
Second Circuit alone could clarify whether its distinctive
environment amplified or moderated the reform's association
with outcomes. Third, merging with the Stanford Securities
Class Action Clearinghouse would provide settlement amount
data to test whether the Second Circuit's low settlement
hazard reflects slower rather than fewer settlements---cases
that take longer to resolve but ultimately settle for larger
amounts. 

\subsection{Methodological Extensions}
Several methodological extensions follow directly from the
limitations documented in Chapter~\ref{ch:discussion}. Full
time-varying coefficient models for circuit, origin, MDL, and
statutory basis would address the proportional hazards
violations that affect all constant hazard-ratio estimates. A
piecewise IPTW model---decomposing the composition-adjusted
PSLRA association across the same three time periods used in
the unweighted analysis---would test whether the time-varying
signature survives composition adjustment. A narrow-window
IPTW analysis restricted to cases filed 1993--1998 would
provide the strongest quasi-experimental test of the PSLRA's
association, though the smaller sample may limit statistical
power. Merging with the Stanford Securities Class Action
Clearinghouse would make settlement amounts available,
enabling a joint model of timing, outcome type, and economic
magnitude.

\subsection{Structural Breaks and Emerging Dynamics}
Testing for post-2018 structural breaks---particularly the
Supreme Court's 2018 decision in \textit{Cyan, Inc.\ v.\
Beaver County Employees Retirement Fund}\footnote{138 S. Ct.\
1061 (2018).}, which expanded state court jurisdiction over
Securities Act claims---is a natural extension of the
piecewise hazard approach developed here. If \textit{Cyan}
redirected Section~11 claims from federal to state courts,
the composition of the federal securities docket may have
shifted in ways the competing-risks framework is well
positioned to detect. Similarly, the rapid growth of SPAC-
and cryptocurrency-related litigation since 2020 may
introduce a new class of cases with distinctive duration and
outcome patterns that warrant separate analysis within the
competing-risks framework.

\section{Closing Reflection}

The central methodological lesson of this thesis is that the
competing-risks framework materially changes the empirical
picture of securities litigation dynamics. A standard
single-outcome analysis of the same data would have found a
significant PSLRA association with dismissal timing, a
significant settlement suppression, and meaningful circuit
variation---but it would have missed the MDL sign reversal,
mischaracterized the settlement association, and provided no
tools for separating the reform's effect from concurrent
compositional shifts. The competing-risks framework does more than fine-tune these estimates; it uncovers structural aspects of the litigation process that a standard approach cannot identify.

More broadly, this thesis suggests that empirical legal
scholarship stands to gain from the same methodological tools
that have transformed biomedical and actuarial research. The
IDB is one of the most comprehensive longitudinal records of
federal civil litigation ever assembled, yet it has been
underutilized by researchers with formal training in survival
analysis. The same competing-risks structure---multiple
mutually exclusive outcomes, large administrative datasets,
and the need to separate time-to-event from event-type---
applies directly to bankruptcy proceedings, employment
discrimination cases, and patent disputes.

The federal securities docket does not resolve neatly. Cases
settle, are dismissed, linger in MDL consolidation, or exit
for reasons the data cannot explain. Any framework that
pretends otherwise---by treating one outcome as the only
outcome, or by averaging over a time-varying process---will
mischaracterize the system it purports to study. A
competing-risks approach takes that complexity seriously, and
in doing so, comes closer to the truth.
